% Imporditud: tools/import_eksperimendid.py
% Allikas: Eesti füüsikaolümpiaadi eksperimendi ülesannete
%   kogu (Rael Kalda, 2023), ülesanne Ü155, lahendus L155.
\setAuthor{Erkki Tempel}
\setRound{piirkonnavoor}
\setYear{2020}
\setNumber{G E1}
\setDifficulty{4}
\setTopic{staatika}
\setVahendid{plastikust kohvitops, A4 paber, joonlaud}
\setPunktid{10}
\setAllikas{Eksperimendi kogu Ü155, lahendus lk 116}
\setLahendusseis{toores_ilma_skeemita}

\prob{Masskese}
Määrake plastiktopsi masskeskme kõrgus lauapinnast.

\hint

\solu
Kasutame paberilehte kaldpinnana ning leiame, millise nurga korral kukub tops ümber. Selleks, et tops ei hakkaks paberil libisema, tuleb paberi serv murda tagasi nii, et kohvitops jääb paberi serva taha pidama. Selleks, et paber ei painduks läbi, saab kasutada mitmekordset paberilehte. Kohvitops tuleb asetada paberilehele nii, et tassi sang asetseks külje peal, siis jääb masskese kaldpinna suhtes tassi keskele.

M

B \_ \_ AO

C

Kolmnurgast $\triangle$ABC saame avaldada

AC tan $\alpha$ =

BC Kolmnurgast $\triangle$AM O saame leida masskeskme M kõrguse kohvitopsi põhjast M O.

tan $\alpha$ = AO $\Rightarrow$ MO = AO $\cdot$ BC

MO AC

Mõõdame pikkused AO, AC ja BC ning arvutame välja masskeskme kõrguse lauapinnast.

\textit{Žürii hindamisskeem on kogumikus lk 116; siia ei ole seda ümber kirjutatud.}
\probend
